\section{Finite moment inversion and recovery}
\label{sec:moment-recovery}

We isolate the elementary inverse principle used below.  It is stated over
$\CC$, although all applications have positive rational bases.

\begin{lemma}[Finite exponential moments]
\label{lem:finite-moments}
Let $z_1,\ldots,z_r$ be distinct nonzero complex numbers and
\[
 u_m=\sum_{i=1}^r a_i z_i^m\qquad(m\geq1),
\]
where $a_i\neq0$.  The sequence $(u_m)_{m\geq1}$ determines the unordered
collection $\{(z_i,a_i):1\leq i\leq r\}$.  If the bases $z_i$ are already
known and $u_m=\sum_i a_i z_i^m$ is available at any $r$ consecutive
nonnegative indices $m=m_0,\ldots,m_0+r-1$, then those moments determine all
coefficients, allowing $m_0=0$ and allowing some coefficients to be zero.
\end{lemma}

\begin{proof}
As a formal power series at the origin,
\begin{equation}
 \sum_{m\geq1}u_m t^{m-1}
 =\sum_{i=1}^r\frac{a_i z_i}{1-z_i t}.
 \label{eq:moment-generating}
\end{equation}
The reduced rational function on the right has simple poles $z_i^{-1}$.
Their positions recover the bases, and their residues recover the
coefficients.  Hence equality of all moments forces equality of the unordered
collections.

When the bases are known, moments with nonnegative indices
$m_0,\ldots,m_0+r-1$ give a
linear system with matrix
$(z_i^{m_0+j})_{0\leq j<r,1\leq i\leq r}$.  Its determinant is a nonzero
monomial times the Vandermonde product
$\prod_{i<j}(z_j-z_i)$, so all coefficients are uniquely determined.
\end{proof}

For each degree $d$ occurring in $\Irr(\K)$, introduce the multiplicities
\begin{align*}
 c_d^+&=\#\{\chi:d_\chi=d,\ \nu_\chi=+1\},\\
 c_d^-&=\#\{\chi:d_\chi=d,\ \nu_\chi=-1\},\\
 c_d^0&=\#\{\chi:d_\chi=d,\ \nu_\chi=0\},\\
 t_d&=c_d^++c_d^-+c_d^0.
\end{align*}

\begin{theorem}[Reconstruction from the two spectra]
\label{thm:reconstruction}
The joint sequences
\[
 (O_{\K}(m))_{m\geq1},\qquad (N_{\K}(n))_{n\geq1}
\]
determine $|\K|$ and every multiplicity $c_d^+,c_d^-,c_d^0$.  Equivalently,
they determine the multiset
$\{(d_\chi,\nu_\chi):\chi\in\Irr(\K)\}$.
\end{theorem}

\begin{proof}
We proceed in four steps.

\emph{Step 1: the group order.}
By \eqref{eq:orientable-spectrum},
\[
 O_{\K}(m)^{1/(4m)}
 =|\K|\left(\sum_{\chi\in\Irr(\K)}d_\chi^{-2m}\right)^{1/(4m)}.
\]
There is at least one degree-one character, namely the trivial character.  The
sum in parentheses is bounded below by $1$ and above by
$|\Irr(\K)|$, independently of $m$.  Therefore
\begin{equation}
 |\K|=\lim_{m\to\infty}O_{\K}(m)^{1/(4m)}.
 \label{eq:recover-order}
\end{equation}

\emph{Step 2: all degree multiplicities.}
After recovering $|\K|$, normalize the orientable moments as
\begin{equation}
 P_m=\frac{O_{\K}(m)}{|\K|^{4m}}
     =\sum_d t_d(d^{-2})^m.
 \label{eq:P-moments}
\end{equation}
Inverse-degree sums $\sum_{\chi}\chi(1)^{-t}$ are standard character-degree
zeta values; see, for example, Liebeck and Shalev \cite{LiebeckShalev2005}.
We use this standard finite-group expression at positive even integers, not a new zeta
function.  Distinct degrees give distinct nonzero bases $d^{-2}$, and every
coefficient $t_d$ is positive.  By \cref{lem:finite-moments}, $(P_m)$ recovers
the bases and coefficients, hence every occurring degree $d$ and its total
multiplicity $t_d$.

\emph{Step 3: the self-dual multiplicities.}
At even nonorientable indices, indicator-zero characters vanish and both
nonzero indicators have even power.  Thus
\begin{equation}
 Q_m=\frac{N_{\K}(2m)}{|\K|^{4m}}
 =\sum_d(c_d^++c_d^-)(d^{-2})^m,
 \qquad m\geq1.
 \label{eq:Q-moments}
\end{equation}
The bases are already known from \eqref{eq:P-moments}.  The Vandermonde part
of \cref{lem:finite-moments}, with zero coefficients allowed, recovers
\[
 s_d:=c_d^++c_d^-
\]
for every occurring $d$.

\emph{Step 4: the indicator signs.}
At odd index $n=2m+1$, where $m\geq0$,
\begin{align}
 R_m&=\frac{N_{\K}(2m+1)}{|\K|^{4m+2}}\notag\\
 &=\sum_d(c_d^+-c_d^-)d^{-(2m+1)}
  =\sum_d\frac{c_d^+-c_d^-}{d}(d^{-2})^m.
 \label{eq:R-moments}
\end{align}
Again the bases are known.  If $r$ distinct degrees occur, the known-base
clause of \cref{lem:finite-moments}, now with $m_0=0$, shows that
$R_0,\ldots,R_{r-1}$ recover the coefficients
\[
 b_d:=\frac{c_d^+-c_d^-}{d}.
\]
Since each degree $d$ is already known from \eqref{eq:P-moments}, multiplying
by $d$ gives $\delta_d:=d b_d=c_d^+-c_d^-$ for every degree.  Finally,
\begin{equation}
 c_d^+=\frac{s_d+\delta_d}{2},\qquad
 c_d^-=\frac{s_d-\delta_d}{2},\qquad
 c_d^0=t_d-s_d.
 \label{eq:multiplicity-recovery}
\end{equation}
This reconstructs the claimed multiset.
\end{proof}

\begin{corollary}[Exact equivalence of the selected periodic data]
\label{cor:spectral-equivalence}
For finite groups $\K$ and $\K'$, the two identities
\[
 O_{\K}(m)=O_{\K'}(m)\quad(m\geq1),\qquad
 N_{\K}(n)=N_{\K'}(n)\quad(n\geq1)
\]
hold if and only if
\[
 |\K|=|\K'|
 \quad\text{and}\quad
 \{(d_\chi,\nu_\chi):\chi\in\Irr(\K)\}
 =\{(d_\psi,\nu_\psi):\psi\in\Irr(\K')\}
\]
as multisets.
\end{corollary}

\begin{proof}
The forward implication is \cref{thm:reconstruction}.  For the reverse
implication, substitute the common order and pair multiset into
\eqref{eq:orientable-spectrum} and \eqref{eq:nonorientable-spectrum}.
\end{proof}

\begin{remark}[Finite algebra after the order limit]
\label{rem:finite-recovery}
Once $|\K|$ is known and there are $r$ distinct degrees, the remaining
reconstruction is finite.  The rational function
\eqref{eq:moment-generating} for $(P_m)$ has denominator degree $r$; standard
finite recurrence recovery determines it from finitely many exact moments.
After its $r$ bases are known, $r$ even nonorientable moments and $r$ odd
nonorientable moments solve two Vandermonde systems.  The infinite sequences
in \cref{thm:reconstruction} are used to state the invariant without assuming
$r$ or $|\K|$ in advance.
\end{remark}

The parity split is essential.  The orientable moments alone cannot distinguish
characters of equal degree with different indicators.  The even
nonorientable moments detect self-duality but cannot distinguish $+1$ from
$-1$.  The odd nonorientable moments supply exactly the missing signed data.
